\chapter{Discrete Random Variables}

\begin{center}\itshape
Shuffle a deck of $52$ cards and follow $1\to\sigma(1)\to\sigma(\sigma(1))\to\cdots$ until you come back to $1$.\\
The deck splits into cycles. How many cycles do you expect: about $5$, about $15$, or about $26$?
\end{center}

A random variable is a number that depends on chance, and its expectation is the average value it would take
over many repetitions. The one tool of this chapter is \emph{linearity of expectation}: the expectation of a sum is
the sum of the expectations, whether or not the terms are independent. Written as a sum of indicators, a
complicated count often has an expectation that one can compute in one line, even when its law is out of reach.

The chapter sets up discrete random variables, their expectation and variance, and the classical laws
(binomial, geometric, Poisson), then applies the indicator trick to three gems: the cycles of a random
permutation, the convergence of fixed points to the Poisson law promised in Volume~1, and the coupon collector.

\section{Random variables and their laws}

\begin{definition}
A \emph{discrete random variable} on a probability space $(\Omega,\mathcal F,P)$ is a map $X:\Omega\to E$, where
$E$ is a countable set, such that $\{X=x\}=\{\omega:X(\omega)=x\}$ is an event for every $x\in E$. Its \emph{law} is
the probability $\mu_X$ on $E$ given by $\mu_X(\{x\})=P(X=x)$. Discrete random variables $X_1,\dots,X_n$ (with values
in $E_1,\dots,E_n$) are \emph{independent} if
\[
P(X_1=x_1,\dots,X_n=x_n)=\prod_{i=1}^nP(X_i=x_i)
\]
for all $x_i\in E_i$; an infinite family is independent if every finite subfamily is.
\end{definition}

The events $\{X=x\}$ for $x\in E$ partition $\Omega$, so $\sum_xP(X=x)=1$ and $\mu_X$ is a probability on $E$ by
Proposition~\ref{prop:discrete}. For $A\subseteq E$, $\{X\in A\}$ is the countable union of the $\{X=x\}$, $x\in A$,
hence an event, of probability $\mu_X(A)$. A pair $(X,Y)$ of discrete random variables is again one, with values in
the countable set $E\times F$; its law is the \emph{joint law}. Real-valued discrete random variables are those with
$E\subseteq\R$; sums and products of them are again discrete random variables (their values form a countable set).

\begin{proposition}\label{prop:indepfun}
If $X_1,\dots,X_n$ are independent, then for all sets $A_i\subseteq E_i$, $P(X_1\in A_1,\dots,X_n\in A_n)=\prod_iP(X_i\in
A_i)$. Consequently, for $1\le k<n$ and maps $f$, $g$, the random variables $f(X_1,\dots,X_k)$ and $g(X_{k+1},\dots,X_n)$
are independent.
\end{proposition}
\begin{proof}
The event $\{X_1\in A_1,\dots,X_n\in A_n\}$ is the disjoint union of the events $\{X_1=x_1,\dots,X_n=x_n\}$ over
$(x_1,\dots,x_n)\in A_1\times\dots\times A_n$, a countable set. So its probability is
$\sum P(X_1=x_1)\cdots P(X_n=x_n)$, and a sum of non-negative products over a product of countable sets is the
product of the sums (sum by packets, Volume~1, Chapter~3). For the second claim, let $U=(X_1,\dots,X_k)$ and
$V=(X_{k+1},\dots,X_n)$. For tuples $u=(u_1,\dots,u_k)$ and $v=(v_{k+1},\dots,v_n)$, independence gives
$P(U=u,V=v)=\prod_{i\le k}P(X_i=u_i)\prod_{i>k}P(X_i=v_i)=P(U=u)P(V=v)$, where the last step is the first claim with
singletons for some indices and the whole $E_i$ for the others. Summing over $u\in f^{-1}(a)$ and $v\in g^{-1}(b)$ (packets
again) gives $P(f(U)=a,g(V)=b)=P(f(U)=a)P(g(V)=b)$.
\end{proof}

\section{Expectation}

Expectation is a weighted sum, and Volume~1 (Chapter~3) showed that infinite sums behave like finite ones exactly
when they converge absolutely. This is why the definition asks for absolute convergence.

\begin{definition}
Let $X$ be a real discrete random variable. If $X\ge0$, its \emph{expectation} is $\mathbb E[X]=\sum_xx\,P(X=x)\in[0,+\infty]$.
In general, $X$ is \emph{integrable} if $\mathbb E|X|=\sum_x|x|P(X=x)<\infty$, and then
$\mathbb E[X]=\sum_xx\,P(X=x)$, an absolutely convergent series.
\end{definition}

\begin{proposition}[transfer formula]\label{prop:transfer}
Let $X$ be a discrete random variable with values in $E$ and $f:E\to\R$. Then $f(X)$ is integrable iff
$\sum_{x\in E}|f(x)|P(X=x)<\infty$, and then $\mathbb E[f(X)]=\sum_{x\in E}f(x)P(X=x)$. The same holds for $f\ge0$
without condition, in $[0,+\infty]$.
\end{proposition}
\begin{proof}
For each value $y$ of $f(X)$, $P(f(X)=y)=\sum_{x:f(x)=y}P(X=x)$. So $\sum_x|f(x)|P(X=x)$, grouped in packets according
to the value $y=f(x)$, is $\sum_y|y|P(f(X)=y)$: the two series of non-negative terms have the same sum. When it is
finite, the same grouping is allowed for the absolutely convergent series $\sum_xf(x)P(X=x)$, and gives
$\sum_yyP(f(X)=y)=\mathbb E[f(X)]$.
\end{proof}

\begin{theorem}\label{thm:linearity}
Let $X$, $Y$ be real discrete random variables and $a,b\in\R$.
\begin{enumerate}[label=(\alph*)]
\item \emph{Linearity}: if $X$ and $Y$ are integrable, so is $aX+bY$, and $\mathbb E[aX+bY]=a\mathbb E[X]+b\mathbb E[Y]$.
The same holds for $X,Y\ge0$ and $a,b\ge0$, in $[0,+\infty]$.
\item \emph{Positivity}: if $X\le Y$ and both are integrable, $\mathbb E[X]\le\mathbb E[Y]$. For an event $A$, $\mathbb E[\mathbf
1_A]=P(A)$.
\item \emph{Independence}: if $X$ and $Y$ are integrable and independent, $XY$ is integrable and
$\mathbb E[XY]=\mathbb E[X]\mathbb E[Y]$.
\item \emph{Tail sum}: if $N$ takes values in $\N$, $\mathbb E[N]=\sum_{k\ge1}P(N\ge k)$.
\end{enumerate}
\end{theorem}
\begin{proof}
(a) Apply the transfer formula to the pair $Z=(X,Y)$ and $f(x,y)=ax+by$. Integrability:
$\sum_{x,y}|ax+by|P(X=x,Y=y)\le|a|\sum_{x,y}|x|P(X=x,Y=y)+|b|\sum_{x,y}|y|P(X=x,Y=y)$, and summing first over $y$
(packets again), $\sum_{x,y}|x|P(X=x,Y=y)=\sum_x|x|P(X=x)$. The same computation without absolute values, now
legitimate, gives the formula. For non-negative variables, every term is non-negative and no condition is needed.

(b) $Y-X\ge0$ has a non-negative expectation, and use (a). $\mathbf 1_A$ takes the value $1$ with probability $P(A)$.

(c) By transfer with $f(x,y)=|xy|$ and independence,
$\mathbb E|XY|=\sum_{x,y}|x||y|P(X=x)P(Y=y)=\mathbb E|X|\,\mathbb E|Y|<\infty$; the same computation without absolute values gives
the formula.

(d) $N=\sum_{k\ge1}\mathbf 1_{N\ge k}$ (exactly $N$ terms equal $1$). So $\sum_nnP(N=n)=\sum_n\sum_{k=1}^nP(N=n)=
\sum_k\sum_{n\ge k}P(N=n)=\sum_kP(N\ge k)$, exchanging two sums of non-negative terms.
\end{proof}

Linearity needs no independence at all, and this is what makes it powerful. Volume~1 used it to show that a random
permutation has on average one fixed point: the count is a sum of $n$ indicators of probability $\frac1n$. Section~2.5
pushes the idea further.

\section{Variance and covariance}

\begin{proposition}[Cauchy--Schwarz]\label{prop:cs}
If $\mathbb E[X^2]$ and $\mathbb E[Y^2]$ are finite, then $XY$ is integrable and $\bigl(\mathbb E|XY|\bigr)^2\le\mathbb E[X^2]\,\mathbb E[Y^2]$. In
particular $X$ is integrable, and $(\mathbb E|X|)^2\le\mathbb E[X^2]$.
\end{proposition}
\begin{proof}
$|xy|\le\frac12(x^2+y^2)$, so $XY$ is integrable. For real $t$, $0\le\mathbb E[(|X|+t|Y|)^2]=\mathbb E[X^2]+2t\,\mathbb E|XY|+t^2\mathbb E[Y^2]$. If
$\mathbb E[Y^2]=0$ then $Y=0$ with probability $1$ and both sides vanish; otherwise this quadratic in $t$ has no two
distinct real roots, so its discriminant is $\le0$: $(\mathbb E|XY|)^2\le\mathbb E[X^2]\mathbb E[Y^2]$. Take $Y=1$ for the last claim.
\end{proof}

\begin{definition}
For $X$, $Y$ with finite second moments, the \emph{covariance} is $\operatorname{Cov}(X,Y)=\mathbb E[(X-\mathbb EX)(Y-\mathbb EY)]=\mathbb E[XY]-\mathbb E[X]\mathbb E[Y]$, and
the \emph{variance} is $\operatorname{Var}X=\operatorname{Cov}(X,X)=\mathbb E[X^2]-\mathbb E[X]^2$. Its square root $\sigma_X$ is the \emph{standard deviation}.
\end{definition}

\begin{proposition}\label{prop:varsum}
For $X_1,\dots,X_n$ with finite second moments,
\[
\operatorname{Var}\Bigl(\sum_{i=1}^nX_i\Bigr)=\sum_{i=1}^n\operatorname{Var}X_i+2\sum_{i<j}\operatorname{Cov}(X_i,X_j),
\]
and $|\operatorname{Cov}(X_i,X_j)|\le\sigma_{X_i}\sigma_{X_j}$. If the $X_i$ are pairwise independent, the variance of the sum is the sum of the
variances.
\end{proposition}
\begin{proof}
With $Y_i=X_i-\mathbb EX_i$, the variance of the sum is $\mathbb E[(\sum Y_i)^2]=\sum_{i,j}\mathbb E[Y_iY_j]$ by linearity, which is the formula. The
bound is Cauchy--Schwarz applied to $Y_i$, $Y_j$. If $X_i$ and $X_j$ are independent, so are $Y_i$ and $Y_j$
(Proposition~\ref{prop:indepfun}), and $\mathbb E[Y_iY_j]=\mathbb E[Y_i]\mathbb E[Y_j]=0$.
\end{proof}

Only \emph{pairwise} independence is used: this is why the pairwise independent bits of Exercise~\ref{ex:pairwise}
are enough for many randomised algorithms, whose analysis only needs variances.

\section{The classical laws}

\begin{definition}
Let $p\in\,]0,1[$, $n\ge1$ and $\lambda>0$.
\begin{itemize}
\item $X$ is \emph{Bernoulli}$(p)$ if $P(X=1)=p$ and $P(X=0)=1-p$.
\item $X$ is \emph{binomial}$(n,p)$ if $P(X=k)=\binom nkp^k(1-p)^{n-k}$ for $0\le k\le n$.
\item $T$ is \emph{geometric}$(p)$ if $P(T=k)=(1-p)^{k-1}p$ for $k\ge1$.
\item $X$ is \emph{Poisson}$(\lambda)$ if $P(X=k)=e^{-\lambda}\frac{\lambda^k}{k!}$ for $k\ge0$.
\end{itemize}
\end{definition}

These are probabilities on $\N$ by the binomial formula, the geometric series and the exponential series. They
describe, respectively: one trial with success probability $p$; the number of successes in $n$ independent trials;
the index of the first success in a sequence of independent trials; and, as Theorem~\ref{thm:poissonlimit} shows,
the number of rare events among many independent opportunities.

\begin{proposition}\label{prop:laws}
\begin{enumerate}[label=(\alph*)]
\item If $X_1,\dots,X_n$ are independent Bernoulli$(p)$, then $X_1+\dots+X_n$ is binomial$(n,p)$.
\item Means and variances: Bernoulli $p$ and $p(1-p)$; binomial $np$ and $np(1-p)$; geometric $\frac1p$ and $\frac{1-p}{p^2}$;
Poisson $\lambda$ and $\lambda$.
\item The geometric law is \emph{memoryless}: $P(T>j+k\mid T>j)=P(T>k)=(1-p)^k$ for $j,k\ge0$.
\end{enumerate}
\end{proposition}
\begin{proof}
(a) The sum equals $k$ iff exactly $k$ of the $X_i$ equal $1$; each of the $\binom nk$ choices of the set of indices has
probability $p^k(1-p)^{n-k}$ by independence, and these events are disjoint.

(b) Bernoulli: $\mathbb E X=\mathbb E X^2=p$. Binomial: by (a), linearity and Proposition~\ref{prop:varsum}, $np$ and $np(1-p)$ (the law
of the sum, hence its mean and variance, does not depend on how it is realised). Geometric: by the tail sum,
$\mathbb E T=\sum_{k\ge1}P(T\ge k)=\sum_{k\ge1}(1-p)^{k-1}=\frac1p$, since $T\ge k$ iff the first $k-1$ trials fail. Similarly
$T^2=\sum_{k=1}^T(2k-1)$, so $\mathbb E[T^2]=\sum_{k\ge1}(2k-1)P(T\ge k)=2\sum_{k\ge1}kx^{k-1}-\frac1p$ with $x=1-p$. Exchanging two sums
of non-negative terms, $\sum_{k\ge1}kx^{k-1}=\sum_{k\ge1}\sum_{j\ge k}x^{j-1}=\sum_{k\ge1}\frac{x^{k-1}}{1-x}=\frac1{(1-x)^2}=\frac1{p^2}$. So
$\operatorname{Var}T=\frac2{p^2}-\frac1p-\frac1{p^2}=\frac{1-p}{p^2}$. Poisson:
$\mathbb E X=\sum_{k\ge1}e^{-\lambda}\frac{\lambda^k}{(k-1)!}=\lambda$ and similarly $\mathbb E[X(X-1)]=\lambda^2$, so $\operatorname{Var}X=\lambda^2+\lambda-\lambda^2=\lambda$.

(c) $P(T>k)=\sum_{i>k}(1-p)^{i-1}p=(1-p)^k$, so $P(T>j+k)/P(T>j)=(1-p)^k$.
\end{proof}

\begin{theorem}[Poisson's law of rare events]\label{thm:poissonlimit}
If $p_n\in[0,1]$ and $np_n\to\lambda>0$, then for every $k\ge0$, $\binom nkp_n^k(1-p_n)^{n-k}\to e^{-\lambda}\frac{\lambda^k}{k!}$.
\end{theorem}
\begin{proof}
Write
\[
\begin{aligned}\binom nkp_n^k(1-p_n)^{n-k}&=\frac{n(n-1)\cdots(n-k+1)}{n^k}\eqbr\cdot\frac{(np_n)^k}{k!}(1-p_n)^{n}(1-p_n)^{-k}.\end{aligned}
\]
For fixed $k$, the first factor tends to $1$, the second to $\frac{\lambda^k}{k!}$, and the last to $1$ since $p_n\to0$. For
the third, $\ln(1-p)=-p+O(p^2)$ as $p\to0$, so $n\ln(1-p_n)=-np_n+O(np_n^2)\to-\lambda$, because $np_n^2=(np_n)p_n\to0$.
\end{proof}

\begin{figure}[ht]
\centering
\fig{ch02-poisson}
\caption{The binomial laws with $n=6$ and $n=30$ and mean $3$, against the Poisson law of parameter $3$.}
\label{fig:poisson}
\end{figure}

The number of misprints on a page, of calls to a switchboard in a minute, or of soldiers of a Prussian cavalry
corps killed by horse kicks in a year (Bortkiewicz's famous data of 1898) follow Poisson laws for this reason:
many independent opportunities, each with a tiny probability. Figure~\ref{fig:poisson} shows how fast the
approximation becomes good.

\section{Cycles and fixed points of a random permutation}

\begin{think}
In the question that opens the chapter, the number of cycles of a random permutation of $52$ cards is a sum of
indicators, but of which events? The cycle containing a given card is a natural object: how long is it, on
average?
\end{think}

\begin{figure}[ht]
\centering
\fig{ch02-cycles}
\caption{A uniformly random permutation of $20$ elements, drawn cycle by cycle: here $4$ cycles, of lengths $10$,
$8$, $1$ and $1$. The expected number of cycles is $H_{20}\approx3.6$.}
\label{fig:cycles}
\end{figure}

Write a permutation in \emph{canonical cycle notation}: each cycle starts with its smallest element, and the cycles
are listed in increasing order of their first elements, for example
\[
\begin{aligned}&(1\,5\,12\,8\,2\,16\,7\,4)\eqbr(3\,6\,14\,13\,9\,17\,19\,18\,20\,15)(10)(11).\end{aligned}
\]
The notation can be written from left to right by $n$ successive choices, which is the key to the next theorem.

\begin{theorem}[Feller's coupling]\label{thm:feller}
Let $C_n$ be the number of cycles of a uniformly random permutation of $\{1,\dots,n\}$. Then $C_n$ has the law of
$B_1+B_2+\dots+B_n$, where the $B_j$ are independent and $B_j$ is Bernoulli$(\frac1j)$. In particular
\[
\begin{gathered}\mathbb E[C_n]=H_n=\sum_{j=1}^n\frac1j,\qq\operatorname{Var}C_n=\sum_{j=1}^n\frac1j\Bigl(1-\frac1j\Bigr)=H_n-\sum_{j=1}^n\frac1{j^2}.\end{gathered}
\]
\end{theorem}
\begin{proof}
Write the canonical notation as follows. Start with ``$(1$''. At step $k=1,\dots,n$, let $r$ be the number of elements
not yet written; choose either to close the current cycle, or to write one of the $r$ unused elements next. After a
closing, if unused elements remain, open a new cycle with the smallest of them (this is forced, not a choice).

Before step $k$ there are exactly $n-k$ unused elements: before step $1$ there are $n-1$, and each step uses one
element, either by writing it or, for a closing followed by an opening, by the forced opening. So step $k$ offers
$n-k+1$ choices, and the last step forces the final closing. A sequence of choices produces a permutation in canonical
notation: each opening uses the smallest unused element, so every cycle starts with its smallest element and the
first elements increase. Conversely, the canonical notation of a permutation determines the sequence of choices
that writes it. So we have a bijection between the $n\cdot(n-1)\cdots1=n!$ sequences of choices and the permutations.

Make the choices independently, step $k$ uniformly among its $n-k+1$ options. Each sequence of choices then has
probability $\frac1{n!}$, so the permutation produced is uniform. The number of cycles is the number of closings,
$C_n=\sum_{k=1}^nI_k$, where $I_k$ indicates a closing at step $k$; the $I_k$ are functions of independent choices,
hence independent, and $P(I_k=1)=\frac1{n-k+1}$. With $j=n-k+1$, this is the claim, and the mean and variance follow
from Propositions~\ref{prop:laws} and~\ref{prop:varsum}.
\end{proof}

For $52$ cards, $H_{52}\approx4.54$: a shuffled deck has about four or five cycles, with a standard deviation of
$\sqrt{H_{52}-\sum1/j^2}\approx1.71$. Since $H_n=\ln n+\gamma+o(1)$ (Volume~1, Chapter~3), a random permutation of a
million elements has only about $14$ cycles. Most elements lie in a few long cycles (Exercise~\ref{ex:prisoners}).

The fixed points tell a different story: their number converges in law, without any normalisation.

\begin{theorem}\label{thm:fixedpoisson}
Let $F_n$ be the number of fixed points of a uniformly random permutation of $\{1,\dots,n\}$. For every $j\ge0$,
$P(F_n=j)\to e^{-1}\frac1{j!}$ as $n\to\infty$: $F_n$ converges in law to the Poisson law of parameter $1$.
\end{theorem}
\begin{proof}
A permutation with exactly $j$ fixed points is determined by the set of its fixed points, $\binom nj$ choices, and a
permutation without fixed points of the $n-j$ other elements, $D_{n-j}$ choices. With the formula
$\frac{D_m}{m!}=\sum_{i=0}^m\frac{(-1)^i}{i!}$ of Volume~1 (Proposition~2.3),
\[
\begin{aligned}P(F_n=j)&=\binom nj\frac{D_{n-j}}{n!}=\frac1{j!}\cdot\frac{D_{n-j}}{(n-j)!}\eqbr=\frac1{j!}\sum_{i=0}^{n-j}\frac{(-1)^i}{i!}\longrightarrow\frac1{j!}e^{-1}.\end{aligned}
\]
\end{proof}

This is the law of rare events again, in disguise: each element is a fixed point with the small probability $\frac1n$,
and the $n$ events are almost independent. Exercise~\ref{ex:factorial} shows how close to Poisson the law of $F_n$ is,
through its factorial moments.

\section{The coupon collector}

\begin{think}
Each box of cereal contains one of $50$ different toys, uniformly at random. How many boxes do you expect to buy
to collect all $50$? Guess: $50$, $100$, $200$, or $1000$?
\end{think}

Draw coupons $U_1,U_2,\ldots$ independently and uniformly among $n$ types, and let $T_n$ be the first time all
types have appeared. (A probability space carrying an infinite sequence of independent draws is admitted, as in
Chapter~1; all the events below only involve finitely many draws.) For $1\le i\le n$, let $G_i$ be the number of
draws needed, after the $(i-1)$-th new type has appeared, to see the $i$-th new type. Then $T_n=G_1+\dots+G_n$.

\begin{theorem}\label{thm:coupon}
The $G_i$ are independent and $G_i$ is geometric with parameter $p_i=\frac{n-i+1}n$. Consequently
\[
\begin{gathered}\mathbb E[T_n]=nH_n,\qquad\operatorname{Var}T_n\le\frac{\pi^2}6n^2,\qq P(T_n>n\ln n+cn)\le e^{-c}\quad(c>0).\end{gathered}
\]
\end{theorem}
\begin{proof}
Fix $k_1,\dots,k_n\ge1$ and let $K=k_1+\dots+k_n$. The event $\{G_1=k_1,\dots,G_n=k_n\}$ depends on the first $K$ draws,
and we count the sequences of $K$ draws in it. During phase $i$ (the $k_i$ draws after the $(i-1)$-th new type),
the first $k_i-1$ draws must repeat one of the $i-1$ types already seen, and the last one must be one of the
$n-i+1$ new types. So there are $\prod_i(i-1)^{k_i-1}(n-i+1)$ such sequences among $n^K$. Dividing by $n^K$ and noting that
$\frac{i-1}n=1-p_i$ and $\frac{n-i+1}n=p_i$,
\[
P(G_1=k_1,\dots,G_n=k_n)=\prod_{i=1}^n(1-p_i)^{k_i-1}p_i.
\]
This product form shows that the $G_i$ are independent and geometric (sum over all $k_j$, $j\ne i$, to get the law
of $G_i$). Then $\mathbb E[T_n]=\sum_i\frac n{n-i+1}=nH_n$, and $\operatorname{Var}T_n=\sum_i\frac{1-p_i}{p_i^2}\le\sum_i\frac{n^2}{(n-i+1)^2}\le
n^2\sum_{j\ge1}\frac1{j^2}=\frac{\pi^2}6n^2$ (Volume~1, Exercise~6.7).

For the tail, $T_n>m$ iff some type is missing from the first $m$ draws. A given type is missing with probability
$(1-\frac1n)^m\le e^{-m/n}$, so by the union bound $P(T_n>m)\le ne^{-m/n}$, which is $e^{-c}$ for $m=n\ln n+cn$.
\end{proof}

\begin{figure}[ht]
\centering
\fig{ch02-coupon}
\caption{The law of the number $T$ of draws needed to collect $n=50$ coupons, computed exactly. Its mean is
$50H_{50}\approx225$, and $P(T>n\ln n+2n)\approx0.12$, below the bound $e^{-2}\approx0.135$ of Theorem~\ref{thm:coupon}.}
\label{fig:coupon}
\end{figure}

For $50$ toys, $nH_n\approx225$ boxes (Figure~\ref{fig:coupon}); the last toy alone takes $50$ boxes on average, as
many as the first $25$ together. The collector's law is lopsided: it has a long tail to the right, because the last
few types are slow to appear.

\begin{beauty}
Each coupon is a type that must be seen once, and the coupon collector governs every situation where one waits for
all of $n$ equally likely outcomes. A random walk on a complete graph covers it after about $n\ln n$ steps;
randomised load balancing leaves a server idle until about $n\ln n$ requests have arrived; and Chapter~8 shows that
a deck of $n$ cards needs about $n\ln n$ random swaps before it looks shuffled, because until then some card has never
been touched, which is a coupon still missing.
\end{beauty}

\section{Exercises}

\begin{exercise}[$\star$]\label{ex:records}
In a uniformly random order of $n$ distinct numbers, the $k$-th number is a \emph{record} if it is larger than all
the numbers before it. Show that the $k$-th is a record with probability $\frac1k$, that the record events are
independent, and that the number of records has mean $H_n$.
\end{exercise}

\begin{exercise}[$\star$]\label{ex:fixedvar}
Show that the number of fixed points of a uniformly random permutation of $n\ge2$ elements has variance exactly $1$,
using indicators and covariances.
\end{exercise}

\begin{exercise}[$\star\star$]\label{ex:quicksort}
\emph{Randomised quicksort.} To sort $n$ distinct numbers $x_{(1)}<\dots<x_{(n)}$, quicksort picks a uniformly random
pivot, compares it with every other element, splits them into smaller and larger, and sorts both parts recursively
with independent fresh pivots. (a) Show that $x_{(i)}$ and $x_{(j)}$, $i<j$, are compared iff the first pivot chosen
among $x_{(i)},\dots,x_{(j)}$ is one of the two, and that this has probability $\frac2{j-i+1}$. (b) Deduce that the expected
number of comparisons is $2(n+1)H_n-4n\sim2n\ln n$.
\end{exercise}

\begin{exercise}[$\star\star$]\label{ex:factorial}
For a random variable $X$ with values in $\N$, the \emph{factorial moments} are $\mathbb E[(X)_k]$, where
$(x)_k=x(x-1)\cdots(x-k+1)$. (a) Show that the Poisson law of parameter $\lambda$ has $\mathbb E[(X)_k]=\lambda^k$. (b) Show that the
number $F_n$ of fixed points of a random permutation of $n$ elements has $\mathbb E[(F_n)_k]=1$ for $k\le n$: its first $n$
factorial moments are exactly those of the Poisson law of parameter $1$.
\end{exercise}

\begin{exercise}[$\star\star$]\label{ex:reservoir}
\emph{Reservoir sampling.} A stream $x_1,x_2,\dots,x_n$ of unknown length passes by, and we want a uniformly random
sample of $k$ of its elements while storing only $k$. Keep the first $k$ elements; then, for $i>k$, with probability
$\frac ki$ replace a uniformly chosen stored element by $x_i$, independently of the past. Show by induction on $n$ that
after $n$ elements the stored set is a uniformly random $k$-subset of $\{x_1,\dots,x_n\}$.
\end{exercise}

\begin{exercise}[$\star\star$]\label{ex:prisoners}
(a) Show that the number of permutations of $\{1,\dots,m\}$ having a cycle of a given length $k>\frac m2$ is
$\frac{m!}k$, and deduce that a random permutation of $2n$ elements has a cycle longer than $n$ with probability
$H_{2n}-H_n\to\ln2$. (b) \emph{The hundred prisoners.} $100$ prisoners each have a number; their numbers are placed at random in $100$
boxes. Each prisoner may open $50$ boxes, without communicating, and all are freed if each finds their own number.
Show that if prisoner $i$ opens box $i$, then the box whose label is the number found, and so on, they all succeed
with probability $1-(H_{100}-H_{50})\approx0.31$.
\end{exercise}

\begin{exercise}[$\star\star\star$]\label{ex:birthdaywait}
Throw balls uniformly into $N$ boxes until the first collision, and let $W_N$ be the number of balls thrown. (a)
Show that $\mathbb E[W_N]=\sum_{n=0}^Nq_n$, with $q_n$ as in Proposition~\ref{prop:birthday}. (b) Show that
$\mathbb E[W_N]\sim\sqrt{\pi N/2}$ as $N\to\infty$, using the Gaussian integral $\int_0^\infty e^{-x^2/2}\dd x=\sqrt{\pi/2}$ (Volume~1,
Exercise~6.7).
\end{exercise}

\begin{exercise}[$\star\star$]\label{pb:leastqnr}
Let $p$ be an odd prime and $n_p$ the least positive quadratic non-residue modulo $p$.
\begin{enumerate}[label=(\alph*)]
\item Show that $n_p$ is a prime and that $n_p<\sqrt p+1$.
\item Pretend that the symbols $\bigl(\frac qp\bigr)$, $q=2,3,5,7,\dots$ prime, are independent fair signs. Show that the expected
value of $n_p$ in this model is $\sum_{k\ge1}p_k2^{-k}$, where $p_k$ is the $k$-th prime, and compute it to three decimals.
\end{enumerate}
\end{exercise}

\begin{exercise}[$\star\star$]\label{pb:glfix}
Let $A$ be a uniformly random matrix in $GL_n(\mathbb F_q)$, $n\ge2$, and $d=\dim\ker(A-I)$.
\begin{enumerate}[label=(\alph*)]
\item Show that $\mathbb E[q^d]=2$ and $\mathbb E[q^{2d}]=q+3$.
\item Deduce that
\[
\frac1q\le P(1\text{ is an eigenvalue of }A)\le\frac1{q-1}.
\]
\end{enumerate}
\end{exercise}

\begin{exercise}[$\star$]\label{pb:foata}
Write a permutation $\sigma$ in cycle notation with each cycle starting with its largest element, and the cycles in increasing
order of their first elements; then erase the parentheses. For example $(3\,1)(5\,2\,4)\mapsto3\,1\,5\,2\,4$.
\begin{enumerate}[label=(\alph*)]
\item Show that this is a bijection from $\mathfrak S_n$ to the words that use each of $1,\dots,n$ once, and that the number of
cycles of $\sigma$ is the number of records of the word (letters larger than all the letters before them).
\item Deduce that the number of records of a random word and the number of cycles of a random permutation have the same
law, and that exactly $\stirlingi nk$ words have $k$ records.
\end{enumerate}
\end{exercise}

\begin{exercise}[$\star\star$]\label{pb:paley}
\emph{Paley graphs.} Let $p\equiv1\pmod4$ be prime. The Paley graph $P_p$ has vertex set $\mathbb F_p$, with $x$ and $y$ joined
when $x-y$ is a non-zero square (this is symmetric because $-1$ is a square).
\begin{enumerate}[label=(\alph*)]
\item Show that two adjacent vertices have $\frac{p-5}4$ common neighbours and two non-adjacent vertices have $\frac{p-1}4$.
Deduce that $P_p$ has $\frac{p(p-1)(p-5)}{48}$ triangles.
\item For $S,T\subset\mathbb F_p$ let $e(S,T)$ be the number of pairs $(s,t)\in S\times T$ with $s,t$ adjacent. Show that
\[
\Bigl|e(S,T)-\frac{|S||T|-|S\cap T|}2\Bigr|\le\frac12\sqrt{p\,|S|\,|T|}.
\]
\item Show that $P_p$ has no clique and no independent set with more than $\sqrt p$ vertices.
\item Deduce from $P_{17}$ that the Ramsey number $R(4,4)$ is at least $18$.
\end{enumerate}
\end{exercise}

\begin{exercise}[$\star\star$]\label{pb:goodman}
\emph{Goodman's bound (1959).} Colour the edges of $K_n$ red and blue.
\begin{enumerate}[label=(\alph*)]
\item Count the triangles that are not monochromatic through their two-coloured ``corners'', and show that there are at least
\[
\binom n3-\frac n2\Bigl\lfloor\frac{(n-1)^2}4\Bigr\rfloor
\]
monochromatic triangles. Check that this gives $2$ for $n=6$.
\item Show that this is asymptotically $\frac14\binom n3$, and that a random colouring has $\frac14\binom n3$ monochromatic triangles on average.
\end{enumerate}
\end{exercise}

\begin{exercise}[$\star\star$]\label{pb:influence}
\emph{Influences.} For a Boolean $f$, the \emph{influence} of voter $i$ is $\operatorname{Inf}_i(f)=P\bigl(f(x)\ne f(x^{\oplus i})\bigr)$, the probability that $i$ is
pivotal, and the \emph{total influence} is $I(f)=\sum_i\operatorname{Inf}_i(f)$.
\begin{enumerate}[label=(\alph*)]
\item Show that $\operatorname{Inf}_i(f)=\sum_{S\ni i}\hat f(S)^2$ and $I(f)=\sum_S|S|\hat f(S)^2$. Deduce the Poincaré inequality $\operatorname{Var}f\le I(f)$, with
equality if and only if $f$ has degree at most $1$.
\item Call $f$ \emph{monotone} if changing a vote from $-1$ to $1$ never changes the outcome from $1$ to $-1$. Show that then
$\operatorname{Inf}_i(f)=\hat f(\{i\})$, and deduce that $I(f)\le\sqrt n$ for every monotone $f$.
\item Show that for $n$ odd, $\operatorname{Inf}_i(\operatorname{Maj}_n)=\binom{n-1}{(n-1)/2}2^{1-n}$, that $I(\operatorname{Maj}_n)\sim\sqrt{\frac{2n}\pi}$, and that $W^1(\operatorname{Maj}_n)\to\frac2\pi$.
\end{enumerate}
\end{exercise}

\begin{exercise}[$\star\star\star$]\label{pb:arrow}
\emph{Condorcet, Arrow, and Kalai's formula.} Three candidates $a$, $b$, $c$; each of $n$ voters ranks them, choosing one of the six
rankings uniformly at random, independently of the others. A Boolean $f$ decides each of the three duels: let $x_i=1$ if voter $i$
prefers $a$ to $b$ and $x_i=-1$ otherwise, and similarly $y_i$ for $b$ against $c$ and $z_i$ for $c$ against $a$; society prefers $a$ to
$b$ if $f(x)=1$, and so on. The outcome is \emph{rational} if it is not a cycle, that is, unless $f(x)=f(y)=f(z)$.
\begin{enumerate}[label=(\alph*)]
\item Show that each $(x_i,y_i,z_i)$ is uniform on the six vectors of $\{-1,1\}^3$ other than $\pm(1,1,1)$, and deduce that each of the
pairs $(x,y)$, $(y,z)$, $(z,x)$ is $(-\frac13)$-correlated.
\item (Kalai, 2002.) Show that
\[
P(\text{rational})=\frac34-\frac34\operatorname{Stab}_{-1/3}(f).
\]
\item Show that for majority with $n=3$ this probability is $\frac{17}{18}$, and that for $\operatorname{Maj}_n$ it tends to
$\frac34+\frac3{2\pi}\arcsin\frac13=0.9123\ldots$ (Guilbaud, 1952).
\item (Arrow's theorem for three candidates.) Suppose that the outcome is rational for every profile, and that $f$ is
\emph{unanimous}: $f(1,\dots,1)=1$. Show that $f$ is a dictatorship, $f(x)=x_i$ for some $i$.
\end{enumerate}
\end{exercise}

\begin{exercise}[$\star\star\star$]\label{pb:lcs}
\emph{Common subsequences of random words.} Let $A$ and $B$ be independent uniform words of length $n$ over $\{0,1\}$, and $L_n$ the length of their longest common
subsequence (not necessarily contiguous).
\begin{enumerate}[label=(\alph*)]
\item Show that $\mathbb E[L_{m+n}]\ge\mathbb E[L_m]+\mathbb E[L_n]$, and deduce that $\frac{\mathbb E[L_n]}n$ converges to $\gamma_2=\sup_n\frac{\mathbb E[L_n]}n$.
\item By exact enumeration one finds numerically $\mathbb E[L_8]=\frac{44803}{8192}$. What lower bound for $\gamma_2$ follows?
\item Show that $P(L_n\ge k)\le\binom nk^22^{-k}$. Let $H(\kappa)=-\kappa\ln\kappa-(1-\kappa)\ln(1-\kappa)$; using $\binom nk\le e^{nH(k/n)}$, show that $\gamma_2\le\kappa^*$, where $\kappa^*$ is the root in
$\left]\frac12,1\right[$ of $2H(\kappa)=\kappa\ln2$. Compute $\kappa^*$.
\end{enumerate}
\end{exercise}

\begin{figure}[ht]
\centering
\fig{ch10-condorcet}
\caption{The probability that majority rule gives a rational outcome with three candidates and $n$ random voters, computed
exactly from Kalai's formula (Problem~\ref{pb:arrow}), for $n=1,3,\dots,51$. It starts at $1$ (one voter is a dictator), drops to
$\frac{17}{18}$ for three voters, and decreases to Guilbaud's limit $0.9123$: about one election in eleven has no Condorcet winner.}
\label{fig:condorcet}
\end{figure}

\section{Hints}

\begin{hint}{ex:records}
The relative order of the first $k$ numbers is uniform, and does not depend on the relative order of the first
$k-1$.
\end{hint}
\begin{hint}{ex:fixedvar}
For $i\ne j$, $P(\sigma(i)=i,\sigma(j)=j)=\frac1{n(n-1)}$.
\end{hint}
\begin{hint}{ex:quicksort}
(a) Until a pivot falls among $x_{(i)},\dots,x_{(j)}$, these elements stay in the same part. (b) Group the pairs by $d=j-i$.
\end{hint}
\begin{hint}{ex:factorial}
(b) $(F_n)_k$ counts the ordered $k$-tuples of distinct fixed points.
\end{hint}
\begin{hint}{ex:reservoir}
Show that each $k$-subset of the first $n$ elements has probability $1/\binom nk$; split according to whether $x_{n}$ is in it.
\end{hint}
\begin{hint}{ex:prisoners}
(a) Choose the elements of the long cycle, arrange them in a cycle, permute the rest; there is at most one cycle
longer than $\frac m2$. (b) Prisoner $i$ follows the cycle of the permutation that contains $i$.
\end{hint}
\begin{hint}{ex:birthdaywait}
(a) Tail sum. (b) Bound $q_n$ above by $e^{-n(n-1)/2N}$ and, for $n\le N^{3/5}$, below with $\ln(1-x)\ge-x-x^2$ on
$[0,\frac12]$; compare the sums with integrals.
\end{hint}

\begin{hint}{pb:leastqnr}
(a) Look at $m=\lfloor p/n_p\rfloor+1$ and $mn_p-p$.
\end{hint}

\begin{hint}{pb:glfix}
$GL_n(\mathbb F_q)$ acts on vectors, and on pairs of vectors; count the orbits. For (b), use $q^d-1\ge q-1$ when $d\ge1$,
and Cauchy--Schwarz.
\end{hint}

\begin{hint}{pb:foata}
To go back from the word, cut it just before each record.
\end{hint}

\begin{hint}{pb:paley}
(a) Count with $\frac{1+\chi}2$ and use Problem~\ref{pb:conics}(a). (b) Cauchy--Schwarz in $s$, then sum over all $s\in\mathbb F_p$.
(c) If $K$ is a clique and $n$ a non-square, $(c,c')\mapsto c+nc'$ is injective on $K\times K$.
\end{hint}

\begin{hint}{pb:goodman}
A triangle that is not monochromatic has exactly two two-coloured corners. A vertex of red degree $r$ is the corner of
$r(n-1-r)$ two-coloured angles.
\end{hint}

\begin{hint}{pb:influence}
(a) $f(x)\ne f(x^{\oplus i})$ iff the derivative $D_if(x)=\frac12\bigl(f(x^{i\to1})-f(x^{i\to-1})\bigr)$ is $\pm1$. (b) For monotone $f$, $D_if\in\{0,1\}$.
Cauchy--Schwarz. (c) Voter $i$ is pivotal iff the others split evenly. Stirling.
\end{hint}

\begin{hint}{pb:arrow}
(a) List the six rankings. (b) For $u,v,w\in\{-1,1\}$, $1-1_{u=v=w}=\frac34-\frac14(uv+vw+wu)$. (c) For odd $f$, $\operatorname{Stab}_{-\rho}(f)=-\operatorname{Stab}_\rho(f)$.
(d) $(-\frac13)^k\ge-\frac13$, with equality only for $k=1$. A Boolean function of degree $1$ is $\pm x_i$.
\end{hint}

\begin{hint}{pb:lcs}
(a) Cut both words into a prefix of length $m$ and a suffix of length $n$. (b) Superadditivity gives $\gamma_2\ge\frac{\mathbb E[L_n]}n$ for every $n$. (c) Choose the $k$ positions in each word;
they match with probability $2^{-k}$. Then $\mathbb E[L_n]\le k+nP(L_n\ge k)$.
\end{hint}

\section{Solutions}

\begin{solution}{ex:records}
\idea Records depend only on relative orders, and the relative order of a random arrangement is uniform at every
scale.
\writeup Let $R_k$ be the indicator that the $k$-th number is a record, i.e.\ the largest of the first $k$. In a uniform
order, the position of the largest among the first $k$ is uniform in $\{1,\dots,k\}$, so $P(R_k=1)=\frac1k$. For
independence, it suffices to see that for $k_1<\dots<k_m$, $P(R_{k_1}=\dots=R_{k_m}=1)=\frac1{k_1\cdots k_m}$ (the events
with some $R$'s equal to $0$ follow by Proposition~\ref{prop:indepcomp}). Build the relative order of the first $k$
numbers by inserting the $k$-th number among the first $k-1$: the $k!$ relative orders correspond bijectively to
the choices of a rank $r_1\in\{1\},r_2\in\{1,2\},\dots,r_k\in\{1,\dots,k\}$ for each new number among those before it, so
these ranks are independent and uniform. $R_k=1$ iff $r_k=k$, and so the $R_k$ are independent, as functions of
independent ranks. Linearity gives $\mathbb E[\sum R_k]=H_n$. (Compare with Theorem~\ref{thm:feller}: the same law $\sum B_j$
appears, which is no accident.)
\end{solution}

\begin{solution}{ex:fixedvar}
\idea Write $F=\sum_i\mathbf 1_{\sigma(i)=i}$ and add up all the covariances.
\writeup Let $I_i=\mathbf 1_{\sigma(i)=i}$. Then $\operatorname{Var}I_i=\frac1n(1-\frac1n)$, and for $i\ne j$, $\operatorname{Cov}(I_i,I_j)=\frac1{n(n-1)}-\frac1{n^2}=\frac1{n^2(n-1)}$.
By Proposition~\ref{prop:varsum}, $\operatorname{Var}F=n\cdot\frac{n-1}{n^2}+n(n-1)\cdot\frac1{n^2(n-1)}=\frac{n-1}n+\frac1n=1$. The indicators are slightly
positively correlated, and the correlation exactly compensates the $\frac1n$ missing from $n$ independent Bernoulli$(\frac1n)$.
\end{solution}

\begin{solution}{ex:quicksort}
\idea Two elements are compared only when one of them is the pivot and the other is still in the same part; follow
the block $x_{(i)},\dots,x_{(j)}$.
\writeup (a) As long as no pivot has been chosen in the block $B=\{x_{(i)},\dots,x_{(j)}\}$, every pivot is outside $B$, hence
smaller or larger than all of $B$, and $B$ stays inside one part. Let $y$ be the first pivot chosen in $B$. If $y$ is
$x_{(i)}$ or $x_{(j)}$, it is compared with all of its part, including the other one. Otherwise $x_{(i)}<y<x_{(j)}$ are sent to
different parts and are never compared (comparisons only happen with a pivot). By symmetry, given that the pivot
is the first chosen in $B$, it is uniform in $B$: the probability is $\frac2{j-i+1}$. (Formally, the first time a pivot is
drawn from a part containing $B$, conditional on landing in $B$ it is uniform on $B$.)
(b) By linearity, the expected number of comparisons is $\sum_{i<j}\frac2{j-i+1}=\sum_{d=1}^{n-1}(n-d)\frac2{d+1}$. With $m=d+1$,
this is $\sum_{m=2}^n\frac{2(n+1-m)}m=2(n+1)(H_n-1)-2(n-1)=2(n+1)H_n-4n$. Since $H_n\sim\ln n$, it is $\sim2n\ln n$.
\end{solution}

\begin{solution}{ex:factorial}
\idea Factorial moments count ordered tuples, which is exactly what fixed points make easy.
\writeup (a) $\mathbb E[(X)_k]=\sum_{j\ge k}j(j-1)\cdots(j-k+1)e^{-\lambda}\frac{\lambda^j}{j!}=\lambda^k\sum_{j\ge k}e^{-\lambda}\frac{\lambda^{j-k}}{(j-k)!}=\lambda^k$.
(b) $(F_n)_k$ is the number of ordered $k$-tuples $(i_1,\dots,i_k)$ of distinct fixed points, so by linearity
$\mathbb E[(F_n)_k]=\sum P(\sigma(i_1)=i_1,\dots,\sigma(i_k)=i_k)$ over the $n(n-1)\cdots(n-k+1)$ ordered tuples. Each term is
$\frac{(n-k)!}{n!}$ (the other $n-k$ elements are permuted freely), and the product is $1$.
\end{solution}

\begin{solution}{ex:reservoir}
\idea Track the probability of each final $k$-subset, and peel off the last element of the stream.
\writeup Claim: after $n\ge k$ elements, each $k$-subset $S$ of $\{x_1,\dots,x_n\}$ is the stored set with probability
$1/\binom nk$. For $n=k$ the stored set is all of them. Assume it for $n-1$. If $x_n\notin S$, then $S$ was stored at time
$n-1$ (probability $1/\binom{n-1}k$) and $x_n$ was rejected (probability $1-\frac kn$): total $\frac{n-k}n/\binom{n-1}k=1/\binom nk$. If
$x_n\in S$, then at time $n-1$ the stored set was $S\setminus\{x_n\}$ plus one element $y$ among the $n-k$ others, which
was then replaced by $x_n$: $(n-k)\cdot\frac1{\binom{n-1}k}\cdot\frac kn\cdot\frac1k=\frac{n-k}{n\binom{n-1}k}=1/\binom nk$.
\end{solution}

\begin{solution}{ex:prisoners}
\idea Long cycles are rare enough to count directly, because two of them cannot coexist. In (b), the strategy
turns the boxes into a permutation, and every prisoner follows their own cycle.
\writeup (a) Choose the $k$ elements of the cycle ($\binom mk$ ways), arrange them in a cycle ($(k-1)!$ ways) and permute
the others ($(m-k)!$ ways): $\binom mk(k-1)!(m-k)!=\frac{m!}k$ permutations. Two cycles longer than $\frac m2$ would need more
than $m$ elements, so these sets for different $k>\frac m2$ are disjoint, and for $m=2n$ the probability of a cycle
longer than $n$ is $\sum_{k=n+1}^{2n}\frac1k=H_{2n}-H_n$, which tends to $\ln2$ since $H_n-\ln n\to\gamma$.
(b) Let $\sigma(b)$ be the number in box $b$. Prisoner $i$ opens $i,\sigma(i),\sigma^2(i),\dots$: they walk along the cycle of
$\sigma$ containing $i$, and find their number $i$ just before returning to box $i$, after as many boxes as the length of
that cycle. So everyone succeeds iff all cycles have length at most $50$, which has probability
$1-(H_{100}-H_{50})\approx1-0.688=0.312$. Each prisoner alone succeeds with probability $\frac12$, but the successes are
strongly correlated, which is what the strategy exploits.
\end{solution}

\begin{solution}{ex:birthdaywait}
\idea The tail sum turns the mean waiting time into $\sum q_n$, and $q_n$ looks like a Gaussian $e^{-n^2/2N}$ in $n$,
whose sum is close to its integral.
\writeup (a) $W_N>n$ iff the first $n$ balls are in different boxes, so $P(W_N>n)=q_n$, which vanishes for $n>N$. The tail
sum (Theorem~\ref{thm:linearity}(d)) gives $\mathbb E[W_N]=\sum_{n\ge0}P(W_N>n)=\sum_{n=0}^Nq_n$.

(b) Upper bound. $q_n\le e^{-n(n-1)/2N}$ (Proposition~\ref{prop:birthday}) and $n(n-1)=(n-\frac12)^2-\frac14$, so
$q_n\le e^{1/8N}f(n-\frac12)$ with $f(x)=e^{-x^2/2N}$, decreasing on $[0,\infty[$. For $n\ge2$, $f(n-\frac12)\le\int_{n-3/2}^{n-1/2}f$, so
$\sum_{n\ge1}f(n-\frac12)\le f(\frac12)+\int_{1/2}^\infty f\le1+\sqrt N\sqrt{\pi/2}$ (substitute $x=\sqrt Nu$). So
$\mathbb E[W_N]\le1+e^{1/8N}(1+\sqrt{\pi N/2})$.

Lower bound. For $0\le x\le\frac12$, $\ln(1-x)\ge-x-x^2$ (the function $\ln(1-x)+x+x^2$ vanishes at $0$ and has derivative
$\frac{x(1-2x)}{1-x}\ge0$). Let $M=\lfloor N^{3/5}\rfloor$ and $N\ge32$, so that $\frac kN\le\frac12$ for $k<M$. For $n\le M$,
\[
\begin{aligned}\ln q_n&=\sum_{k<n}\ln\Bigl(1-\frac kN\Bigr)\ge-\frac{n^2}{2N}-\frac{n^3}{3N^2}\eqbr\ge-\frac{n^2}{2N}-\frac{N^{-1/5}}3,\end{aligned}
\]
using $\sum_{k<n}k\le\frac{n^2}2$ and $\sum_{k<n}k^2\le\frac{n^3}3$. So
\[
\begin{aligned}\mathbb E[W_N]&\ge e^{-N^{-1/5}/3}\sum_{n=1}^Mf(n)\eqbr\ge e^{-N^{-1/5}/3}\Bigl(\int_0^Mf-1\Bigr),\end{aligned}
\]
and $\int_0^Mf=\sqrt N\int_0^{M/\sqrt N}e^{-u^2/2}\dd u$ with $M/\sqrt N\to\infty$. Dividing both bounds by $\sqrt{\pi N/2}$, they tend to
$1$.
\end{solution}

\begin{solution}{pb:leastqnr}
\idea If $n=n_p$ were composite, its factors, being smaller, would be residues, and so would $n$. For the size, use multiples:
the first multiple of $n$ beyond $p$ is $mn$, and $mn-p<n$ is a residue, so $m$ must be a non-residue, hence $m\ge n$. That forces
$n^2\lesssim p$. In the random model, $n_p=p_k$ when the first $k-1$ prime symbols are $+1$ and the $k$-th is $-1$.

\writeup (a) Every $1\le a<n_p$ is a residue. If $n_p=ab$ with $1<a,b<n_p$, then $\bigl(\frac{n_p}p\bigr)=\bigl(\frac ap\bigr)\bigl(\frac bp\bigr)=1$,
absurd; so $n_p$ is prime. Let $n=n_p$ and $m=\lfloor\frac pn\rfloor+1$, so $0<mn-p<n$ (here $n\nmid p$, because $1<n<p$: the
number $1$ is a residue, and non-residues exist in $\{1,\dots,p-1\}$). Then $mn-p$ is a residue, so
$\bigl(\frac mp\bigr)\bigl(\frac np\bigr)=\bigl(\frac{mn}p\bigr)=\bigl(\frac{mn-p}p\bigr)=1$, $\bigl(\frac mp\bigr)=-1$, and $m\ge n$. So $n\le\frac pn+1$,
$n^2-n\le p$, and $n\le\frac{1+\sqrt{1+4p}}2<\sqrt p+1$.

(b) In the model, $P(n_p=p_k)=2^{-(k-1)}\cdot\frac12=2^{-k}$, so $E[n_p]=\sum_kp_k2^{-k}=\frac22+\frac34+\frac58+\frac7{16}+\dots=3.6746\ldots$.
\tools Volume~2, Chapter~6 (Legendre symbol); Volume~4, Chapter~2 (expectation of a discrete variable).
\leads Erdős (1961) proved that the average of $n_p$ over primes $p\le x$ tends to $\sum p_k2^{-k}$, so the model is right on average; the
one finds numerically $3.61$ for $p<2\cdot10^5$. The worst case is another story: Burgess (1957) showed $n_p\ll p^{1/(4\sqrt e)+\varepsilon}$, still the
best known exponent, while the generalised Riemann hypothesis gives $n_p<2\log^2p$ (Ankeny, Bach), which is what makes Miller's
deterministic test work.
\end{solution}

\begin{solution}{pb:glfix}
\idea Fixed vectors of $A$ form the subspace $\ker(A-I)$, with $q^d$ elements. So $q^d$ is the number of fixed points of $A$ on
$V=\mathbb F_q^n$, and $q^{2d}$ the number on $V\times V$. Count the orbits of $GL_n$ on vectors and on pairs of vectors.

\writeup (a) $GL_n(\mathbb F_q)$ acts on $V$ with two orbits, $\{0\}$ and $V\setminus\{0\}$ (any non-zero vector can be completed
to a basis, and an invertible matrix maps one basis to another). By Burnside $\mathbb E[q^d]=2$. On $V\times V$, the orbit of
$(v,w)$ is determined by the linear relations between $v$ and $w$: the orbits are $\{(0,0)\}$; $\{(0,w):w\ne0\}$; $\{(v,0):v\ne0\}$;
$\{(v,\lambda v):v\ne0\}$ for each $\lambda\in\mathbb F_q^*$; and the pairs of independent vectors, one orbit since $n\ge2$ (complete
$(v,w)$ to a basis). That is $3+(q-1)+1=q+3$ orbits, so $\mathbb E[q^{2d}]=q+3$.

(b) Let $Y=q^d-1\ge0$. Then $\mathbb E[Y]=1$ and $\mathbb E[Y^2]=(q+3)-2\cdot2+1=q$. When $d\ge1$, $Y\ge q-1$, so
$1=\mathbb E[Y]\ge(q-1)P(d\ge1)$. And by Cauchy--Schwarz, $1=\mathbb E[Y\mathbf 1_{d\ge1}]^2\le\mathbb E[Y^2]P(d\ge1)=q\,P(d\ge1)$.
\tools Volume~1, Chapter~7 (linear algebra); Volume~3, Chapter~3 (Burnside); Volume~4, Chapter~2 (Cauchy--Schwarz).
\leads For $q=2$ the bounds are $\frac12$ and $1$, and one finds numerically $0.714$ for $n=3$. The exact law of $d$ comes from the
\emph{cycle index} of $GL_n(\mathbb F_q)$, a $q$-analogue of the cycle index of $\mathfrak S_n$ (see Fulman's survey). As
$n\to\infty$ with $q$ fixed, the probability of the eigenvalue $1$ converges, and the limit tends to $0$ like $\frac1q$.
\end{solution}

\begin{solution}{pb:foata}
\idea In the word, the first letters of the cycles are exactly the records, because each cycle starts with its maximum and
the maxima increase. So the cycles can be read back by cutting before each record.

\writeup (a) In the standard form, the first letter of each cycle is its largest element, and these first letters increase from
cycle to cycle. Each later letter of a cycle is smaller than the cycle's first letter, which is the largest letter seen so far.
So the records of the word are exactly the first letters of the cycles. Conversely, given a word, cut it before each record:
the pieces start with their largest element and their first letters increase, so reading the pieces as cycles gives a
permutation whose standard form is the word. The two maps are inverse to each other, and the number of cycles is the number of
records.

(b) A bijection maps the uniform law to the uniform law, so the two counts have the same law. The number of permutations with
$k$ cycles is $\stirlingi nk$ by definition, and the bijection transports it to the words.
\tools Volume~3, Chapter~4 (cycle decomposition); Volume~4, Chapter~2 (records, Feller's coupling).
\leads This is Foata's \emph{fundamental transformation} (Foata and Schützenberger, 1970). It explains why Volume~4 found the same
law, a sum of independent Bernoulli$(\frac1j)$, for the records (Exercise~2.1) and for the cycles
(Theorem~2.12).
\end{solution}

\begin{solution}{pb:paley}
\idea Adjacency is a value of the quadratic character, so every count in the Paley graph is a quadratic character sum, and those
are all $O(\sqrt p)$ by Problem~\ref{pb:conics}.

\writeup (a) The number of common neighbours of $0$ and $d\ne0$ is
\[
\sum_{z\ne0,d}\frac{1+\chi(z)}2\cdot\frac{1+\chi(z-d)}2=\frac14\Bigl[(p-2)+\sum_{z\ne0,d}\chi(z)+\sum_{z\ne0,d}\chi(z-d)+\sum_z\chi(z^2-dz)\Bigr].
\]
The two middle sums are $-\chi(d)$ and $-\chi(-d)=-\chi(d)$, and the last is $-1$ by Problem~\ref{pb:conics}(a) (discriminant $d^2\ne0$).
So it is $\frac{p-3-2\chi(d)}4$: $\frac{p-5}4$ if $d$ is a square (adjacent), $\frac{p-1}4$ otherwise. By translation the same holds for
any two vertices. Each of the $\frac{p(p-1)}4$ edges lies in $\frac{p-5}4$ triangles and each triangle has $3$ edges, so there are
$\frac{p(p-1)(p-5)}{48}$ triangles.

(b) For $s\ne t$, $[s\sim t]=\frac{1+\chi(s-t)}2$, and $\chi(0)=0$, so
$e(S,T)=\frac{|S||T|-|S\cap T|}2+\frac12\sum_{s\in S,t\in T}\chi(s-t)$. By Cauchy--Schwarz, extending the sum over $s$ to all of
$\mathbb F_p$,
\[
\Bigl|\sum_{s\in S}\sum_{t\in T}\chi(s-t)\Bigr|^2\le|S|\sum_{s\in\mathbb F_p}\Bigl|\sum_{t\in T}\chi(s-t)\Bigr|^2
=|S|\sum_{t,t'\in T}\sum_s\chi\bigl((s-t)(s-t')\bigr).
\]
The inner sum is $p-1$ if $t=t'$ and $-1$ otherwise (Problem~\ref{pb:conics}(a) again), so the total is
$|S|\bigl(|T|(p-1)-|T|(|T|-1)\bigr)=|S||T|(p-|T|)\le p|S||T|$.

(c) Let $K$ be a clique and $n$ a non-square. If $c+nc'=d+nd'$ with $c,c',d,d'\in K$ and $c\ne d$, then $c'\ne d'$, and
$c-d=n(d'-c')$ is a square times a non-square: a non-square, contradicting $c\sim d$. So $c=d$, then $c'=d'$: the map
$(c,c')\mapsto c+nc'$ is injective on $K\times K$, and $|K|^2\le p$. The map $x\mapsto nx$ is an isomorphism from $P_p$ to its
complement, so independent sets are bounded in the same way.

(d) $\sqrt{17}<5$, so by (c) $P_{17}$ has no clique and no independent set of size $5$. A direct check
shows that the largest ones have size $3$. A colouring of the edges of $K_{17}$ by ``square'' and
``non-square'' has no monochromatic $K_4$, so $R(4,4)>17$.

\tools Problem~\ref{pb:conics}; Volume~4, Chapter~2 (Cauchy--Schwarz); Volume~2, Chapter~6.
\leads Greenwood and Gleason (1955) proved $R(4,4)=18$, with $P_{17}$ as the extremal graph. The discrepancy bound (b) is the
definition of a \emph{quasirandom} graph in the sense of Chung, Graham and Wilson (1989): such graphs have about as many copies of
every small subgraph as a random graph. The clique bound $\sqrt p$ is far from the truth, which is believed to be $O(\log^2p)$,
the size of the cliques of a random graph being about $2\log_2p$; improving $\sqrt p$ even to $\sqrt{p/2}$ came only in 2021
(Hanson and Petridis).
\end{solution}

\begin{solution}{pb:goodman}
\idea A triangle that is not monochromatic has two corners where the colours change. Counting such corners vertex by vertex
gives at most $\frac{(n-1)^2}4$ per vertex.

\writeup (a) Call an \emph{angle} a pair of edges at a common vertex, and say it is two-coloured if its edges have different colours. A
monochromatic triangle has no two-coloured angle, and every other triangle has exactly two (with colours $RRB$, the two angles at
the ends of the blue edge). Every angle lies in exactly one triangle. So if $T$ is the number of monochromatic triangles,
\[
2\Bigl(\binom n3-T\Bigr)=\sum_vr_v(n-1-r_v),
\]
where $r_v$ is the red degree of $v$. Since $r(n-1-r)\le\lfloor\frac{(n-1)^2}4\rfloor$ for integers $r$, $T\ge\binom n3-\frac n2\lfloor\frac{(n-1)^2}4\rfloor$.
For $n=6$: $20-3\cdot6=2$.

(b) The bound is $\frac{n^3}6-\frac{n^3}8+O(n^2)=\frac{n^3}{24}+O(n^2)\sim\frac14\binom n3$. In a uniformly random colouring, each triangle is
monochromatic with probability $2\cdot2^{-3}=\frac14$, and linearity of expectation gives $\frac14\binom n3$.

\tools Volume~1, Chapter~2 (double counting); Volume~4, Chapter~2 (linearity of expectation).
\leads Since $R(3,3)=6$ (Volume~4, Appendix~A), every colouring of $K_6$ has a monochromatic triangle; Goodman's bound says
it has at least two, which is sharp. That the random colouring is asymptotically optimal is the first case of a
conjecture of Erdős on $K_q$, which turned out to be false already for $K_4$ (Thomason, 1989).
\end{solution}

\begin{solution}{pb:influence}
\idea Flipping voter $i$ matters exactly where the discrete derivative in the $i$-th direction is non-zero, and the derivative of a
multilinear polynomial is read off its coefficients. Parseval then turns probabilities into sums of squares.

\writeup (a) Let $D_if(x)=\frac12\bigl(f(x^{i\to1})-f(x^{i\to-1})\bigr)$, where $x^{i\to\pm1}$ is $x$ with $x_i$ set to $\pm1$. On the expansion, $D_ix^S=x^{S\setminus\{i\}}$ if
$i\in S$ and $0$ otherwise, so $D_if=\sum_{S\ni i}\hat f(S)x^{S\setminus\{i\}}$. For Boolean $f$, $D_if$ takes values in $\{-1,0,1\}$, and $|D_if(x)|=1$ exactly when
$f(x)\ne f(x^{\oplus i})$. Hence, by Parseval (the monomials $x^{S\setminus\{i\}}$, $S\ni i$, are distinct),
\[
\operatorname{Inf}_i(f)=\mathbb E[(D_if)^2]=\sum_{S\ni i}\hat f(S)^2,\qquad I(f)=\sum_i\sum_{S\ni i}\hat f(S)^2=\sum_S|S|\hat f(S)^2.
\]
And $\operatorname{Var}f=\sum_{S\ne\emptyset}\hat f(S)^2\le\sum_S|S|\hat f(S)^2$, with equality iff $\hat f(S)=0$ whenever $|S|\ge2$.

(b) If $f$ is monotone, $D_if\in\{0,1\}$, so $(D_if)^2=D_if$ and $\operatorname{Inf}_i(f)=\mathbb E[D_if]=\hat f(\{i\})$, the constant term of $D_if$. By Cauchy--Schwarz and
Parseval, $I(f)=\sum_i\hat f(\{i\})\le\sqrt n\bigl(\sum_i\hat f(\{i\})^2\bigr)^{1/2}\le\sqrt n$.

(c) Voter $i$ is pivotal for majority iff the other $n-1$ votes split evenly, which has probability $\binom{n-1}{(n-1)/2}2^{1-n}$. With $n-1=2m$ and
$\binom{2m}m\sim\frac{4^m}{\sqrt{\pi m}}$ (Stirling), $\operatorname{Inf}_i\sim\frac1{\sqrt{\pi m}}\sim\sqrt{\frac2{\pi n}}$, so $I(\operatorname{Maj}_n)=n\operatorname{Inf}_1\sim\sqrt{\frac{2n}\pi}$. Majority is monotone, so
$\hat f(\{i\})=\operatorname{Inf}_i$ by (b), and $W^1(\operatorname{Maj}_n)=n\operatorname{Inf}_1^2\to\frac2\pi$.

\tools{Volume~4, Chapter~2 (binomial law) and Volume~1, Chapter~3 (Stirling's formula); Volume~1, Chapter~1 (Cauchy--Schwarz).}
\leads So among monotone rules, majority gives the voters almost the most power possible: $\sqrt{2/\pi}=0.80$ of the bound $\sqrt n$. At the
other extreme, Kahn, Kalai and Linial (1988) proved that every Boolean $f$ with $\operatorname{Var}f=1$ has a voter with influence at least $c\frac{\log n}n$,
and the ``tribes'' function of Ben-Or and Linial shows that this is sharp; the proof uses the Bonami lemma of Problem~\ref{pb:bonami}.
For $f=1-2\cdot1_A$, $I(f)$ is the number of edges of the cube leaving $A$ divided by $2^{n-1}$, which links (a) with the edge-isoperimetric
inequality of Problem~\ref{pb:cubeiso}.
\end{solution}

\begin{solution}{pb:arrow}
\idea One voter's three pairwise preferences are three correlated signs, and ``not all equal'' is a quadratic polynomial in three
signs. Taking expectations gives noise stability at the correlation $-\frac13$, and the sign of $(-\frac13)^k$ separates level $1$ from all others.

\writeup (a) The six rankings $a{>}b{>}c$, $a{>}c{>}b$, $b{>}a{>}c$, $b{>}c{>}a$, $c{>}a{>}b$, $c{>}b{>}a$ give $(x_i,y_i,z_i)=(1,1,-1)$, $(1,-1,-1)$, $(-1,1,-1)$, $(-1,1,1)$,
$(1,-1,1)$, $(-1,-1,1)$: the six vectors other than $\pm(1,1,1)$, each once (a cycle would need a cyclic individual ranking). Each coordinate is
uniform, and $x_iy_i=1$ for exactly two of the six, so $\mathbb E[x_iy_i]=\frac{2-4}6=-\frac13$; the same holds for $(y_i,z_i)$ and $(z_i,x_i)$ by the cyclic
symmetry $a\to b\to c\to a$. Voters are independent, so the three pairs are $(-\frac13)$-correlated.

(b) For $u,v,w\in\{-1,1\}$: if $u=v=w$, $uv+vw+wu=3$; otherwise exactly one product is $1$ and $uv+vw+wu=-1$. So
$1-1_{u=v=w}=\frac34-\frac14(uv+vw+wu)$. With $u=f(x)$, $v=f(y)$, $w=f(z)$ and (a),
\[
P(\text{rational})=\frac34-\frac14\bigl(\mathbb E[f(x)f(y)]+\mathbb E[f(y)f(z)]+\mathbb E[f(z)f(x)]\bigr)=\frac34-\frac34\operatorname{Stab}_{-1/3}(f).
\]

(c) $\operatorname{Maj}_3=\frac12(x_1+x_2+x_3)-\frac12x_1x_2x_3$ (check the eight points), so $W^1=\frac34$, $W^3=\frac14$, $\operatorname{Stab}_{-1/3}=-\frac14-\frac1{108}=-\frac7{27}$, and
$P=\frac34+\frac7{36}=\frac{17}{18}$: a cycle has probability $\frac1{18}$, that is $12$ of the $216$ profiles, the answer to the think box. Majority is odd,
$f(-x)=-f(x)$, so only odd $|S|$ occur and $\operatorname{Stab}_{-\rho}=-\operatorname{Stab}_\rho$. By Problem~\ref{pb:majoritystable}, $\operatorname{Stab}_{-1/3}(\operatorname{Maj}_n)\to-\frac2\pi\arcsin\frac13$, and
$P\to\frac34+\frac3{2\pi}\arcsin\frac13=0.91226\ldots$; a cycle has limiting probability $0.0877\ldots$ (Figure~\ref{fig:condorcet}).

(d) If every profile is rational, $P=1$ and $\operatorname{Stab}_{-1/3}(f)=-\frac13$, that is $\sum_k\bigl((-\frac13)^k+\frac13\bigr)W^k(f)=0$ since $\sum_kW^k=1$. Each bracket is
$\ge0$, and is $0$ only for $k=1$; so all the spectral weight is on level $1$, and $f=\sum_ia_ix_i$. Flipping $x_i$ changes $f$ by $2a_ix_i\in\{0,\pm2\}$, so
$a_i\in\{0,\pm1\}$, and $\sum a_i^2=1$ leaves one non-zero $a_i$: $f=\pm x_i$. Unanimity excludes $-x_i$.

\tools{Problems~\ref{pb:majoritystable} and~\ref{pb:influence}; Volume~4, Chapter~2 (expectation).}
\leads Arrow's theorem in general (any number of candidates and any rules for the different pairs, under the conditions of
unanimity and independence of irrelevant alternatives) needs a combinatorial proof; Kalai's formula gives the case above and
more importantly, a robust version: if $P(\text{rational})\ge1-\varepsilon$, then by the Friedgut--Kalai--Naor theorem $f$ is $O(\varepsilon)$-close to a
dictatorship or an anti-dictatorship. Mossel (2012) extended this quantitative Arrow theorem to any number of candidates and to different rules for different
pairs.
\end{solution}

\begin{solution}{pb:lcs}
\idea Common subsequences of two halves glue into a common subsequence of the whole, so the mean is superadditive and has a linear growth rate. The first moment
method caps it.

\writeup (a) If $A=A'A''$ and $B=B'B''$ with $|A'|=|B'|=m$, a common subsequence of $(A',B')$ followed by one of $(A'',B'')$ is common to $(A,B)$. The pairs $(A',B')$ and $(A'',B'')$ have the laws of
random pairs of lengths $m$ and $n$, so $\mathbb E[L_{m+n}]\ge\mathbb E[L_m]+\mathbb E[L_n]$. Since $\mathbb E[L_n]\le n$, Fekete's lemma (Problem~\ref{pb:saw}) applied to $-\mathbb E[L_n]$ gives convergence of $\frac{\mathbb E[L_n]}n$ to
$\sup_n\frac{\mathbb E[L_n]}n\le1$.

(b) $\gamma_2\ge\frac{\mathbb E[L_8]}8=\frac{44803}{65536}=0.6836$.

(c) $L_n\ge k$ means that for some index sets $I,J\subset\{1,\dots,n\}$ of size $k$ the subwords $A_I$ and $B_J$ are equal; for fixed $I,J$ they are independent uniform words of length $k$, equal with
probability $2^{-k}$. The union bound gives $P(L_n\ge k)\le\binom nk^22^{-k}$. From $1\ge\binom nkx^k(1-x)^{n-k}$ at $x=\frac kn$ we get $\binom nk\le e^{nH(k/n)}$. The function $\lambda(\kappa)=2H(\kappa)-\kappa\ln2$ is
concave, positive at $\frac12$, equal to $-\ln2$ at $1$, and decreasing on $[\frac12,1]$ (its derivative is $2\ln\frac{1-\kappa}\kappa-\ln2<0$), so it has one root $\kappa^*$ there. For $\kappa>\kappa^*$ and $k=\lceil\kappa n\rceil$,
$P(L_n\ge k)\le e^{n\lambda(\kappa)}$, and $\mathbb E[L_n]\le k+nP(L_n\ge k)\le\kappa n+1+ne^{n\lambda(\kappa)}$. Dividing by $n$, $\gamma_2\le\kappa$; so $\gamma_2\le\kappa^*=0.9051$.

\tools{Volume~4, Chapter~2 (union bound, expectations) and the binomial theorem; Problem~\ref{pb:saw} (Fekete).}
\leads Chv\'atal and Sankoff (1975) proved that $\gamma_2$ exists. Its value is still unknown: Lueker (2009) proved $0.788071\le\gamma_2\le0.826280$, and numerical simulations suggest about
$0.81$. Longest common subsequences compare files (the \texttt{diff} utility) and DNA sequences, where one wants to know how long a common subsequence two unrelated
sequences have by chance.
\end{solution}

\section*{Sources and further reading}

\begin{itemize}
\item W.~Feller, \emph{An Introduction to Probability Theory and Its Applications}, Vol.~1, 3rd ed., Wiley, 1968,
Chapters~IX--XI. Random variables, expectation, the classical laws; Chapter~X.6 for cycles of permutations.
\item R.~Arratia, A.~D.~Barbour and S.~Tavaré, \emph{Logarithmic Combinatorial Structures: a Probabilistic
Approach}, European Mathematical Society, 2003, Chapter~1. Feller's coupling and the cycle structure of random
permutations.
\item L.~von Bortkiewicz, \emph{Das Gesetz der kleinen Zahlen}, Teubner, 1898. The horse-kick data and the law of
rare events.
\item P.~Flajolet and R.~Sedgewick, \emph{Analytic Combinatorics}, Cambridge University Press, 2009 (free online),
Chapter~II. Cycles, records and coupons through generating functions, a theme of Chapter~3.
\item M.~Mitzenmacher and E.~Upfal, \emph{Probability and Computing}, 2nd ed., Cambridge University Press, 2017,
Chapters~2--3. Quicksort, the coupon collector and the variance method.
\item J.~S.~Vitter, \emph{Random sampling with a reservoir}, ACM Transactions on Mathematical Software 11 (1985),
37--57.
\item A.~Gál and P.~B.~Miltersen, \emph{The cell probe complexity of succinct data structures}, ICALP 2003, where the
hundred prisoners problem first appeared; and P.~Winkler, \emph{Mathematical Mind-Benders}, A K Peters, 2007, for
its popular version.
\item P.~Flajolet, D.~Gardy and L.~Thimonier, \emph{Birthday paradox, coupon collectors, caching algorithms and
self-organizing search}, Discrete Applied Mathematics 39 (1992), 207--229. Waiting times for birthdays and coupons.
\end{itemize}
